% HIER BITTE AUSFÜLLEN:

\newcommand{\abschlussarbeit}{Bachelorarbeit} % Bachelorarbeit / Masterarbeit
\newcommand{\zweitpruefer}{Prof. Dr. Erika Mustermann} % Zweitprüfer

\newcommand{\sprache}{english} % Deutsch: german / Englisch: english
\newcommand{\druck}{oneside} %  / Simplex: oneside / Duplex: twoside

\newcommand{\titel}{Titel der Abschlussarbeit\\
                    Titel der Abschlussarbeit}

\newcommand{\name}{Max Mustermann}
\newcommand{\matrikelnummer}{1234567}
\newcommand{\adresse}{Warburger Str. 100, 33098 Paderborn}
\newcommand{\upbmail}{student@mail.upb.de}
\newcommand{\abgabedatum}{01. Januar 2020}


% % % % % % % % % % % % % % % % % % % % % % % % % % % % % % % % % % % % % % %


\documentclass[fontsize=12pt,BCOR=15mm,DIV=15,a4paper,headsepline,headings=small,\druck,openright,appendixprefix]{scrbook}

\usepackage[onehalfspacing]{setspace}
\raggedbottom
\usepackage{emptypage}
\makeatother 
\usepackage{parskip}

\usepackage[utf8]{inputenc} % Zeichenkodierung
\usepackage[T1]{fontenc} % Trennung von Wörtern mit Umlauten
\usepackage{lmodern} % Latin Modern font benutzen (enhanced computer modern clone, serif)
\usepackage[\sprache]{babel} % Dokument in deutsch inklusive Silbentrennung nach neuer deutscher Rechtschreibung
\usepackage{natbib} % Stil für Literaturverzeichnis
\usepackage{amsmath} % verbesserte mathematische Umgebung
\usepackage{amsfonts} % zusätzliche mathematische Zeichen
\usepackage{amssymb} % zusätzliche mathematische Zeichen
\usepackage{amsthm} % zusätzliche mathematische Zeichen
\usepackage[hyphens]{url}
\usepackage{exscale} % Anpassung mathematischer Symbole an die Schriftgröße
\usepackage{amstext} % \text in mathematischer Umgebung
\usepackage{array} % Tabellen
\usepackage{rotating} % rotierte Tabellen und Bilder
\usepackage[algoruled, algochapter]{algorithm2e} % Pseudocode
\usepackage{verbatim} % Funktion zum Schreiben von unformatiertem Text
\usepackage[section]{placeins} % Verhindert das Wandern von floats in eine andere section
\usepackage{color} % Farbiger Text
\usepackage{graphicx} % Funktionen zum Einfügen von Bildern
\usepackage{subfigure} % Anzeigen von Bildern aus mehreren Einzelbildern
\usepackage{setspace} % Einstellen des Zeilenabstandes
\usepackage[urlcolor = black,plainpages=false,pdfpagelabels=true,colorlinks=true,linkcolor=black,citecolor=black,bookmarksopen=true]{hyperref} % Verweise werden im PDF zu Hyperlinks
\usepackage{scrlayer-scrpage} % Koma für Kopfzeilen-Formatierung
\usepackage[nounderscore]{syntax}
\usepackage{tabulary} % bessere Tabellen
\usepackage{tabularx} % bessere Tabellen
\usepackage{longtable} % für Tabellen über mehrere Seiten
\usepackage{array}
\usepackage{paralist} % für kleine Einzüge mit \setdefaultleftmargin{0.5em}{}{}{}{}{}
\usepackage{booktabs}
\usepackage{multirow}
\usepackage[table,xcdraw]{xcolor}
\usepackage[hang]{footmisc} %Große Zahlen in Fußnoten
\makeatletter
\def\@makefnmark{\hbox{\textsuperscript\@thefnmark}}
% \ohead{\leftmark} %Only chapter in header (not section!)
\makeatother
\let\footnotesize\small

\bibliographystyle{apalike} % Literaturverzeichnis entsprechend DIN 1505 formatiert, allerdings ohne URL und ISBN
\renewcommand{\topfraction}{0.85} % Anteil einer Seite, die von Floats belegt sein darf (default 0.7)
\renewcommand{\textfraction}{0.1} % Anteil einer Seite, die mindestens Text sein muss (default 0.2)
\renewcommand{\floatpagefraction}{0.75} % Anteil einer Seite, die ein Float einnehmen darf, ohne auf die nächste Seite gesetzt zu werden (default 0.5)

\usepackage{lipsum}

% Begin
\begin{document}

% Front Matter
\pagestyle{empty}
\pagenumbering{alph}

\vfil

\begin{titlepage}

        \begin{center}
        	\includegraphics[scale=1.]{img/data_analytics_group.png}
        \end{center}
        
        \vspace{8ex}
        
        \begin{center}
            
            Universität Paderborn\\
            Fakultät für Wirtschaftswissenschaften\\
            Department Wirtschaftsinformatik\\
            
            \vspace{8ex}
            
            \Large
            \abschlussarbeit\\
            
            \vspace{7ex}
            
            \textbf{\sffamily{\titel}}
            
            \vspace{7ex}
            
            \normalsize
            
            by\\
            \name\\
            Matrikelnummer:~\matrikelnummer\\
            \adresse\\
            \upbmail\\
            
            \vspace{8ex}
            
            vorgelegt bei\\
            Prof. Dr. Oliver Müller\\
            \zweitpruefer\\
            
            \vspace{10ex}
            
            \abgabedatum
            
        \end{center}
        
    \end{titlepage}
    
\vfil 
\cleardoublepage

\noindent
\section*{Eidesstattliche Erklärung}

Hiermit erkläre ich, \nameFirst, an Eides statt, dass ich die vorliegende Arbeit selbstständig, ohne fremde Hilfe und ohne Benutzung anderer als der angegebenen Hilfsmittel angefertigt habe. Die aus fremden Quellen (einschließlich elektronischer Quellen) direkt oder indirekt übernommenen Gedanken, Tabellen, Skizzen, Zeichnungen, bildliche Darstellungen usw. sind ausnahmslos als solche kenntlich gemacht. Die Arbeit ist in gleicher oder ähnlicher Form oder auszugsweise im Rahmen einer anderen Prüfung noch nicht vorgelegt worden.

\bigskip

Paderborn, den \abgabedatum

\begin{flushright}
\bigskip
\bigskip
\bigskip

------------------------------------------- \\
\nameFirst
\end{flushright}

\section*{Eidesstattliche Erklärung}

Hiermit erkläre ich, \nameSecond, an Eides statt, dass ich die vorliegende Arbeit selbstständig, ohne fremde Hilfe und ohne Benutzung anderer als der angegebenen Hilfsmittel angefertigt habe. Die aus fremden Quellen (einschließlich elektronischer Quellen) direkt oder indirekt übernommenen Gedanken, Tabellen, Skizzen, Zeichnungen, bildliche Darstellungen usw. sind ausnahmslos als solche kenntlich gemacht. Die Arbeit ist in gleicher oder ähnlicher Form oder auszugsweise im Rahmen einer anderen Prüfung noch nicht vorgelegt worden.

\bigskip

Paderborn, den \abgabedatum

\begin{flushright}
\bigskip
\bigskip
\bigskip

------------------------------------------- \\
\nameSecond
\end{flushright}
\cleardoublepage

\section*{Abstrakt}

@bkasten füg mal bitte das Exposé hier ein!

\textbf{Stichworte:} data science, machine learning
\cleardoublepage

\pagestyle{scrheadings}
\clearscrheadfoot
\ohead{\headmark}
\ofoot{\pagemark}

\frontmatter

\tableofcontents
\cleardoublepage

\listoffigures
\cleardoublepage

\listoftables
\cleardoublepage

% Thesis
\mainmatter
\chapter{Introduction}

\lipsum[1-4]

\begin{figure}[h]
\centering
\includegraphics[width=0.5\textwidth]{img/data_analytics_group.png}
\caption[Data Analytics Group]{Data Analytics Group}
\label{fig:data_analytics_group}
\end{figure}

\lipsum[5]
\chapter{Related Work}

\lipsum[6-10]
\chapter{Implementation}

\section{Lorem Ipsum}

\lipsum[11-15]

\section{Different citation styles}

\subsection{Subsection one}

\begin{itemize}
    \item citation in parentheses: \citep{muller2017beyond}\\
    Command: \texttt{\textbackslash citep\{muller2017beyond\}}
    
    \item ... with added text after citation: \citep[][p.15]{muller2017beyond}\\
    Command: \texttt{\textbackslash citep[][p.15]\{muller2017beyond\}}
    
    \item ... with added text before citation: \citep[see][]{muller2017beyond}\\
    Command: \texttt{\textbackslash citep[see][]\{muller2017beyond\}}
    
    \item ... with suppressed parentheses: \citealp{muller2017beyond}\\
    Command: \texttt{\textbackslash citealp\{muller2017beyond\}}
    
    \item ... multiple authors: \citep{muller2017beyond, debertoli2014comparing}\\
    Command: \texttt{\textbackslash citep\{muller2017beyond, debertoli2014comparing\}}
    
    \item author not in parentheses: \citet{muller2017beyond}\\
    Command: \texttt{\textbackslash citet\{muller2017beyond\}}
    
    \item ... with suppressed parentheses: \citealt{muller2017beyond}\\
    Command: \texttt{\textbackslash citealt\{muller2017beyond\}}
    
    \item ... multiple authors: \citet{muller2017beyond, debertoli2014comparing}\\
    Command: \texttt{\textbackslash citet\{muller2017beyond, debertoli2014comparing\}}
    
    \item only author: \citeauthor{muller2017beyond}\\
    Command: \texttt{\textbackslash citeauthor\{muller2017beyond\}}
    
    \item only year: \citeyear{muller2017beyond}\\
    Command: \texttt{\textbackslash citeyear{muller2017beyond}}
    
    \item ... in parentheses: \citeyearpar{muller2017beyond}\\
    Command: \texttt{\textbackslash citeyearpar\{muller2017beyond\}}
\end{itemize}
\chapter{Results}

\lipsum[16]

\begin{table}[h]
\centering
    \begin{tabular}{c c c c c}
        \toprule
        \textbf{Variable} & \textbf{Measurement} & \textbf{Mean} & \textbf{Min.} & \textbf{Max.} \\ 
        \toprule
        Age     & Years         & 26.51     & 17.00     & 40.00     \\
        Height  & Centimeters   & 181.49    & 161.00    & 203.00    \\
        Goals   & total p.s.    & 2.39      & .00       & 50.00     \\
        Assists & total p.s.    & 1.64      & .00       & 20.00     \\
        \bottomrule
    \end{tabular}
\caption[Descriptive Statistics]{Descriptive Statistics \citep{muller2017beyond}}
\label{table:descriptive_statistics}
\end{table}

\lipsum[17-20]
\chapter{Zusammenfassung}

\lipsum[21-25]

% Bibliography
\newpage
\addcontentsline{toc}{chapter}{Bibliography} 
\bibliography{bibliography}

% Appendix (if required...)
% \appendix

% End
\end{document}